\clearpage\begingroup\SourceLeaf{661b}
\ArithmeticHead{26}{L}{H. GEPPERT, ÜBER GAUSS’ ARBEITEN ZUR MECHANIK UND POTENTIALTHEORIE.}
\makeatletter\tagsleft@true\let\veqno\@@leqno\makeatother
\renewcommand{\thefootnote}{\arabic{footnote})}\setcounter{footnote}{45}
\begin{center}\textls[100]{3. Das \textsc{Fourier}sche Ungleichungsprinzip.}\end{center}
\noindent\hspace*{\parindent}Während \textsc{Lagrange} in seiner Mécanique nur Gleichungen zwischen den
Koordinaten bezw. Geschwindigkeiten als Bedingungen für ein System zulässt,
hat \textsc{Fourier}\footnote{\textsc{Fourier}, \textit{Mémoire sur la statique.} Journal de l’École Polytechnique 5, 1798, S. 20; Oeuvres
II, 1890, S. 475—521, insbesondere S. 488.} 1798 erstmalig in einer dem Beweise des Prinzips der virtuellen
Geschwindigkeiten gewidmeten Abhandlung auf die Möglichkeit hingewiesen,
dass die Koordinaten der betrachteten Massenpunkte auch Ungleichungsbedingungen
unterliegen können. Handelt es sich z.\,B. um zwei durch einen masselosen,
unausdehnbaren, biegsamen Faden von der Länge 1 verbundene Punkte,
so können diese sich zwar einander beliebig nähern, aber niemals eine Entfernung,
die grösser als 1 ist, annehmen; hat man eine feste, undurchdringliche Kugel,
so kann zwar ein Massenpunkt auf ihr beliebig gleiten oder rollen, oder sich
von ihrer Oberfläche entfernen, vermag aber nicht in sie einzudringen. Während
für die \textsc{Lagrange}sche Form des Prinzips der virtuellen Verrückungen wesentlich
ist, dass zu jeder Verrückung auch die entgegengesetzte möglich sei, sind
die zuletzt betrachteten Bindungen einseitig, und es gibt virtuelle Verrückungen,
die nicht umkehrbar sind. Das genannte Prinzip gilt dann in der Gleichungsform
nicht mehr, vielmehr tritt an seine Stelle bei \textsc{Fourier} die Bemerkung, dass
das Moment der angreifenden Kräfte bei den erlaubten Verrückungen aus
der Gleichgewichtslage nicht positiv werden darf, und es lautet demnach in
diesem Falle die notwendige und hinreichende Gleichgewichtsbedingung dahin,
dass bei den erlaubten virtuellen Verrückungen
\[
\sum_{i=1}^{3n}X_i\delta x_i\leq0\tag{9}
\]
\noindent sein muss\footnote{Vergl. \textsc{Mach}, a.\,a.\,O. S. 60, wo der Beweis dieses Prinzips gleichzeitig mit dem der virtuellen
Verrückungen entwickelt ist. L. \textsc{Boltzmann}, \textit{Vorlesungen über die Prinzipien der Mechanik} Bd. I, 1897,
S. 141, 216 ff.}.
\clearpage\SourceLeaf{662}\ArithmeticHead{27}{R}{II. DAS PRINZIP DES KLEINSTEN ZWANGES.}
\HistoricalFoot{}{4\textsuperscript{*}}
\textsc{Gauss} hat dieses allgemeine \textsc{Fourier}sche Prinzip immer im Auge gehabt;
er setzt es in der Fassung (9) an die Spitze seiner Abhandlung über die Gleichgewichtsfiguren
der Flüssigkeiten\footnote{Werke V, S. 35.}, und kommt in der eben besprochenen
Abhandlung über das Prinzip des kleinsten Zwanges\footnote{Werke V, S. 27 Anm.} wieder darauf zurück.
Die Einführung von Bedingungsungleichungen neben den von \textsc{Lagrange} allein
betrachteten holonomen oder nichtholonomen Bedingungsgleichungen hat \textsc{Gauss}
stets sehr am Herzen gelegen; auch in dem Briefe an \textsc{Möbius}\footnote{Vergl. Werke XI, 1, S. 17—20.} vom 29. September
1837 kommt er darauf zu sprechen: »Ich habe ein Paarmahl Veranlassung
genommen, das herrliche Prinzip in derjenigen Gestalt, wie ich es
wohl bei mündlichen Vorträgen darzustellen pflege, auszusprechen, z.\,B. im
ersten Artikel meiner \textit{Theoria generalis figurae fluidorum.} Es scheint dies aber
wenig beachtet zu sein, vermuthlich, weil es nur so beiläufig vorkommt, ohne
das Unterscheidende besonders hervorzuheben oder umständlich zu motiviren.
.\, .\, . Wichtiger aber erscheint mir, dass alle Schriftsteller, auch \textsc{Lagrange}, das
Prinzip enger gefasst haben als nöthig ist. Allen Schriftstellern sieht man
die Fessel anhängen, die sie in den engern Raum festbannt, wo nur von
solchen Fällen die Rede ist, die Einer analytischen Formel vollständig unterworfen
sind. Dies aber ist nicht die Erschöpfung der Fälle der Natur. Diese
bietet ebenso oft Fälle dar, wo die Eine analytische Formel nicht ausreicht.«

Die wiederholten Hinweise von \textsc{Gauss} mögen mit Veranlassung dazu gewesen
sein, dass die Frage der Ungleichheitsbedingungen in der Folgezeit
eingehender untersucht wurde.

Zu den holonomen oder nichtholonomen Bedingungsgleichungen (1) mögen
eine Reihe von Ungleichungen der Form
\[
f_\nu(x_1,\, .\, .\, .\, x_{3n},t)\geq0
\qquad\qquad(\nu=1\, .\, .\, .\, l)\tag{10}
\]
\noindent treten. Ist in dem zu untersuchenden Augenblick $t$ für den Bewegungszustand
von $\Sigma$ eine der Funktionen $f_\nu$ positiv, so heisst die entsprechende Bedingung
in diesem Augenblick unwirksam; denn dann sind auch alle hinreichend benachbarten
Bewegungszustände, die den anderen Bedingungen genügen, möglich.
\clearpage\SourceLeaf{663}\ArithmeticHead{28}{L}{H. GEPPERT, ÜBER GAUSS’ ARBEITEN ZUR MECHANIK UND POTENTIALTHEORIE.}
\noindent Verschwindet dagegen $f_\nu$, so müssen die $\dot x_i$ der Bedingung genügen
\[
\sum_{i=1}^{3n}\frac{\partial f_\nu}{\partial x_i}\dot x_i+\frac{\partial f_\nu}{\partial t}\geq0.\tag{11}
\]
\noindent Ist hierin für den durch die $x_i,\dot x_i$ charakterisierten Bewegungszustand die linke
Seite positiv, so ist auch diese Bedingung im Augenblick $t$ unwirksam. Gilt
dagegen in (11) das Gleichheitszeichen, so entsteht für die $\ddot x_i$ die Bedingung
\[
\sum_{i=1}^{3n}\frac{\partial f_\nu}{\partial x_i}\ddot x_i+\varphi_\nu(x_i,\dot x_i,t)\geq0.\tag{12}
\]
\noindent In der Form (11) denken wir uns auch nach passender Abänderung der
Schreibweise die nichtholonomen Ungleichungen einbegriffen, die zu den holonomen
Ungleichungen (10) noch hinzutreten können. Während die Kenntnis
des Bewegungszustandes zu entscheiden gestattet, welche der Bedingungen (10)
oder (11) augenblicklich wirksam ist, wissen wir a priori nicht, welche der
Ungleichungen (12) bei der wirklich stattfindenden Bewegung wirksam sein wird.

Das \textsc{Fourier}sche Prinzip, verbunden mit dem Satze von \textsc{d’Alembert}, besagt
dann, dass für die wirklich eintretende Bewegung
\[
\sum_{i=1}^{3n}(m_i\ddot x_i-X_i)\delta x_i\geq0\tag{13}
\]
\noindent sein muss, wenn die $\delta x_i$ erlaubte Verrückungen bedeuten von der Art, dass auch
die variierte Bahn den wirksamen Ungleichungen (10) und den Bedingungen
(3) genügt. Von \textsc{Gibbs}\footnote{G. W. \textsc{Gibbs}, American Journal of Mathematics, 2, 1879, vgl. auch \textsc{Boltzmann} a.\,a.\,O., Bd. 1,
S. 223—225.} und \textsc{Stäckel}\footnotemark[44] rühren einfache Beispiele her, die
zeigen, dass dieses erweiterte \textsc{Fourier-d’Alembert}sche Prinzip im allgemeinen
nicht zur Bestimmung der Beschleunigungen, also zur eindeutigen Festlegung
der Bewegung ausreicht. Vielmehr zeigt sich gerade in diesem Falle die
grosse Überlegenheit des \textsc{Gauss}schen Satzes.

Spricht man das \textsc{Gauss}sche Prinzip auch hier in der Form aus, dass für die
wirkliche Bahn im Vergleich zu den nach \textsc{Gauss} variierten Bahnen der Zwang
ein Minimum sei, so zeigt sich, dass dieses Prinzip eine wohlbestimmte Bewegung
liefert. Voraussetzung dafür ist nur die Regularität des Systems, nämlich
dass aus den Gleichungen (1) und denjenigen der Beziehungen (11), die
\clearpage\SourceLeaf{664}\ArithmeticHead{29}{R}{II. DAS PRINZIP DES KLEINSTEN ZWANGES.}
\noindent den wirksamen Ungleichungen (12) entsprechen, sich ebenso viele Geschwindigkeitskomponenten
als Funktionen der anderen bestimmen lassen. Nachdem schon
\textsc{Ostrogradsky}\footnote{\textsc{Ph. Ostrogradsky}, \textit{Considérations générales sur les momens des forces.} Mém. de l’Acad. Imp.
d. St. Pétersbourg (6) 1, 1834, S. 565.} sich 1834 mit dem \textsc{Fourier}schen Prinzip befasst hatte, hat
wohl zuerst \textsc{Jacobi} in seinen im W.\,S. 1848/49 gehaltenen Vorlesungen über
Dynamik\footnote{Vergl. \textsc{Voss}, a.\,a.\,O. S. 87, Zitat einer handschriftlichen Ausarbeitung von \textsc{Scheibner}.} das Prinzip des kleinsten Zwanges auf Systeme mit Ungleichheitsbedingungen
angewandt. Ihm folgt \textsc{Ritter} in seiner von \textsc{Gauss} beeinflussten
Dissertation, in der das Prinzip des kleinsten Zwanges auf Systeme mit holonomen
Ungleichungsbedingungen der Form (10) angewandt und die geometrische
Deutung im $3n$-dimensionalen euklidischen Raume erörtert wird, ohne
die rechnerische Durchführung des Ansatzes zu erledigen. Diese mehrdimensionale
Deutung geht sicher auf \textsc{Gauss} zurück, denn wir wissen aus Andeutungen
und Veröffentlichungen, sowie aus dem Bericht von \textsc{Sartorius von
Waltershausen}, dass \textsc{Gauss} sich mehrfach, insbesondere in seinen späteren
Jahren, mit mehrdimensionalen Mannigfaltigkeiten beschäftigt hat.

Fassen wir die $\ddot x_i\sqrt{m_i}$ wieder als Koordinaten eines $R_{3n}$ auf, so bestimmen
die wirksamen Gleichungen der Gesamtzahl $s$ bezw. Ungleichungen der
Form (12) einen Teil eines linearen Unterraumes $R_{3n-s}$, der unter der gemachten
Voraussetzung der Regularität von Hyperebenen begrenzt wird. Die
\textsc{Gauss}sche Aufgabe ist dann damit identisch, vom Punkte mit den Koordinaten
$\dfrac{X_i}{\sqrt{m_i}}$ nach dieser Mannigfaltigkeit den kürzesten Abstand zu ziehen. Der
$R_{3n-s}$ ist überall konvex und einfach-zusammenhängend, wie man aus der
Linearität der Bedingungen (12) leicht ersieht. Daraus folgt, dass es nur ein
Minimum des Abstandes gibt, dass also bei regulärer Lage das \textsc{Gauss}sche
Prinzip den Bewegungszustand eindeutig charakterisiert. Das \textsc{Gauss}sche Prinzip
leistet also im Falle von Ungleichheitsbedingungen mehr als das \textsc{d’Alembert}sche
Prinzip, das es im Besonderen in sich schliesst. \textsc{Gauss} hat daher die
Tragweite seines Grundgesetzes unterschätzt, als er erklärte, es sei mit dem
\textsc{d’Alembert}schen Prinzip identisch. Vielmehr wird man auch hier das \textsc{Gauss}sche
Prinzip als ein neues Axiom auffassen müssen, das im Regularitätsfalle
die wirkliche Bestimmung der Bewegung ermöglicht. Dadurch gewinnt das
\textsc{Gauss}sche Gesetz als einziges den Rang des wahren Grundgesetzes der ana\nobreakdash-
\clearpage\SourceLeaf{665}\ArithmeticHead{30}{L}{H. GEPPERT, ÜBER GAUSS’ ARBEITEN ZUR MECHANIK UND POTENTIALTHEORIE.}
\noindent lytischen Mechanik. Der Grund dafür, dass der \textsc{Gauss}sche Satz sich als umfassender
erweist als der \textsc{d’Alembert}sche, liegt darin, dass der Bereich der
zulässigen Variationen der $\ddot x_i$ im allgemeinen umfassender ist als der der
Variationen der $x_i$.

Will man die angestellten Überlegungen rechnerisch durchführen, so hat
man folgendes Problem zu lösen: es ist das Minimum des Ausdruckes
$z_1^2+ .\, .\, .+z_n^2$ zu suchen, wenn die $z_i$ $r$ linearen Ungleichungen der Form
\[
\varphi_\nu=\sum_{i=1}^{n}a_{\nu i}z_i\geq0
\qquad\qquad(\nu=1\, .\, .\, .\, r)\tag{14}
\]
\noindent unterworfen sind. Gerade mit dieser Aufgabe hat sich \textsc{Gauss} in einer Vorlesung
über die Methode der kleinsten Quadrate im Wintersemester 1850/51 beschäftigt,
die von \textsc{Ritter} ausgearbeitet und in seiner Dissertation verwertet wurde\footnote{Ein Teil der \textsc{Ritter}schen Ausarbeitung ist abgedruckt in Werke X, 1, S. 473—481: der interessierende
Abschnitt aus seiner Dissertation ebendaselbst S. 469—472.}. Der
von \textsc{Gauss} eingeschlagene Weg ist von \textsc{Stäckel}\footnote{P. \textsc{Stäckel}, \textit{Eine von Gauss gestellte Aufgabe des Minimums.} Sitzungsber. d. Heidelb. Akad.
d. Wiss. Abt. A, 1917.} eingehend behandelt worden.
Handelt es sich allgemeiner um die Aufgabe, eine Funktion $f(z_1\, .\, .\, .\, z_n)$ unter
den vorgegebenen Bedingungen (14) zum Minimum zu machen, so deutet man
zunächst wieder die $z_1\, .\, .\, .\, z_n$ als kartesische Koordinaten eines $R_n$, in dem
durch (14) ein Raumstück $\Phi$ ausgesondert wird. Man bestimmt dann zunächst
den Punkt $z_i^0$, der das absolute Minimum von $f$ liefert; liegt er in $\Phi$ oder
auf seiner Begrenzung, so sind die Bedingungen (14) unwirksam, und die Aufgabe
ist durch $z_i^0$ gelöst. Liegt dieser Punkt ausserhalb $\Phi$, so liegt der zu
suchende Punkt des Minimums auf der Begrenzung von $\Phi$. Um ihn zu finden,
verfährt \textsc{Gauss} folgendermassen. Er geht aus von einem Punkte $\zeta_i$ der Begrenzung
von $\Phi$, in dem $n$ der Funktionen $\varphi_\nu$ Null, die übrigen $r-n$ positiv
sind, und bestimmt nun die Richtung, in der man auf der Begrenzung fortschreiten
muss, damit $f$ am schnellsten abnimmt, d.\,h. damit eine gegebene
infinitesimale Änderung von $f$ durch Werte $z_i$ erreicht wird, für die
\[
\sum_{i=1}^{n}(z_i-\zeta_i)^2
\]
\noindent ein Minimum ist. Man erhält so eine Kurve schnellster Abnahme auf der
\clearpage\SourceLeaf{666}\ArithmeticHead{31}{R}{II. DAS PRINZIP DES KLEINSTEN ZWANGES.}
\noindent Begrenzung von $\Phi$, bei deren Durchlaufen man entweder an einen Punkt gelangt,
in dem das Abnehmen von $f$ in Zunehmen übergeht — und dann ist
das ein gesuchter Minimumspunkt, — oder man in eine andere Grenzmannigfaltigkeit
übergeht, auf der man wieder die Kurve schnellster Abnahme
aufzustellen hat. Durch Fortsetzung dieses Verfahrens findet man endlich
die Stelle des Minimums, und diejenigen der Funktionen $\varphi_\nu$, die an dieser
Stelle verschwinden, geben die für das Minimum einzig wirksamen unter den
Ungleichungen (14).

Man kann hiernach die Differentialgleichung der Kurven schnellster Abnahme
aufstellen, auf denen man fortzuschreiten hat. Bei allgemeiner Wahl
von $f$ und $\varphi_\nu$ kann diese Differentialgleichung singuläre Stellen besitzen und
dadurch zur Existenz mehrerer Minima führen, wie dies insbesondere beim
\textsc{Gauss}schen Problem dann eintreten kann, wenn es sich um einen irregulären
Bewegungszustand handelt. Daher bedarf nach Anwendung der \textsc{Gauss}schen
Methode die wirkliche Minimumseigenschaft des ermittelten Punktes noch
eines besonderen Nachweises.

Beim Problem des kleinsten Zwanges ist speziell $f\equiv z_1^2+ .\, .\, .+z_n^2$ zu setzen,
und die $\varphi_\nu$ sind, wie (14) zeigt, linear; dadurch eliminieren sich gerade die
letztgenannten Schwierigkeiten. \textsc{Stäckel} hat gezeigt, dass das \textsc{Gauss}sche Verfahren
hier sogar nach endlich vielen Schritten zu der einzigen Stelle des
Minimums führt, indem der zu durchlaufende Weg schnellster Abnahme jede
Grenzmannigfaltigkeit von $\Phi$ nur endlich oft durchläuft. Die spezielle Wahl
des Ausgangspunktes, die \textsc{Gauss} vornimmt, kann gegebenenfalls unmöglich
sein; jedoch ist diese Wahl unwesentlich, wenn man nur auf der Berandung
von $\Phi$ bleibt. Die unwirksamen Ungleichungen hätte man für die Bestimmung
des Minimums ausscheiden können.

Die Frage, ob durch das Prinzip des kleinsten Zwanges auch im Falle
von Ungleichungen die wirkliche Bewegung eindeutig ermittelt werden kann,
blieb lange unbeantwortet. Auf Anregung von \textsc{Study} untersuchte sie A. \textsc{Mayer}\footnote{A. \textsc{Mayer}, \textit{Über die Aufstellung der Differentialgleichungen der Bewegung für reibungslose Punktsysteme,
die Bedingungsungleichungen unterworfen sind.} Ber. d. Sächs. Akad. d. Wiss. zu Leipzig, Math.-phys.
Classe, 51, 1899, S. 224—244; \textit{Zur Regulierung der Stösse in reibungslosen Punktsystemen, die dem
Zwange von Bedingungsungleichungen unterliegen}, ebenda, S. 245—264.}
und skizzierte ein Verfahren, das vor allem gestattet, die für die wirkliche
\clearpage\SourceLeaf{667}\ArithmeticHead{32}{L}{H. GEPPERT, ÜBER GAUSS’ ARBEITEN ZUR MECHANIK UND POTENTIALTHEORIE.}
\noindent Bewegung wirksamen Bedingungen von den unwirksamen in einem gegebenen
Zeitpunkte zu trennen; für den Fall von ein oder zwei wirksamen Bedingungen
gelang ihm der direkte Nachweis der Eindeutigkeit der Lösung. Der
oben von uns wiedergegebene Eindeutigkeitsbeweis, der hauptsächlich auf der
Linearität der Ungleichungen und der Konvexität des Gebietes beruht, und
der im Wesentlichen wohl auch \textsc{Gauss} nicht entgangen sein dürfte, geht auf
\textsc{Zermelo}\footnote{E. \textsc{Zermelo}, \textit{Über die Bewegung eines Punktsystems bei Bedingungsungleichungen.} Göttinger
Nachrichten, Math.-phys. Klasse 1899, S. 306—310.} und \textsc{Stäckel}\footnotemark[55] zurück. Ähnlich wie im \textsc{Lagrange}schen Falle der
Bedingungsgleichungen lässt sich eine modifizierte Multiplikatorenmethode entwickeln,
die in einer Reihe späterer Arbeiten\footnote{Die vollständige Literatur ist wiedergegeben bei P. \textsc{Stäckel}, \textit{Encykl. d. math. Wiss.} IV, 6,
S. 460 und bei E. B. \textsc{Jourdain}, \textit{Ostwalds Klassiker} Nr. 167, S. 60. Beispiele zur Multiplikatorenmethode
entwickeln \textsc{Mayer} a.\,a.\,O. und \textsc{Henneberg}, Journal f. d. reine u. angew. Math. 113, 1894, S. 179—185.} ausgebaut und auf Beispiele
angewandt worden ist.
\clearpage\endgroup
